% 386_Alpha_Geometrie_De_ch.tex
% Johann Pascher, 1. Okt. 2026
% Wo alpha steckt: Ladungseinheit, zwei Kugeln und die Bezugsenergie

\providecommand{\BM}{\textbf{[B]}}
\providecommand{\KM}{\textbf{[K]}}
\providecommand{\SM}{\textbf{[S]}}
\providecommand{\XM}{\textbf{[X]}}
\providecommand{\QM}{\textbf{[Q]}}

\chapter*{Wo \texorpdfstring{$\alpha$}{alpha} steckt}
\addcontentsline{toc}{chapter}{Wo alpha steckt}

\begin{abstract}
Die FFGFT ist in erster Linie eine verhältnisbasierte Theorie. Sie rechnet in natürlichen Einheiten
mit der Konvention $\alpha=1$ und einer entsprechend umdefinierten Ladung; in dieser Grundform
kommt sie ohne fraktale Korrektur aus. Dieses Dokument zeigt, wohin $\alpha$ bei der
Umdefinition wandert und welche geometrische Bedeutung es dort hat. Mit
$\hbar=c=\varepsilon_0=1$ gilt $e^2=4\pi\alpha$; mit $\alpha=1$ wird $e=\sqrt{4\pi}$, die Ladung
wird über die volle Kugeloberfläche gemessen. $\alpha$ ist dann das Quadrat des Verhältnisses von
historischer zu geometrischer Ladungseinheit und zugleich das Verhältnis zweier Kugeln um das
Elektron, $\alpha=r_e/\lambda_C$. Die FFGFT-Brücke $\alpha=\xi E_0^2$ lautet in T0-Einheiten
$\xi E_0^2=1$, also $E_0^2=1/\xi=7500$. Undurchsichtig ist allein die Bezugsenergie der SI-Form
$\alpha_{\text{SI}}=\xi(E_0/1\,\text{MeV})^2$: Die Übereinstimmung sagt, dass diese Bezugsenergie
auf $10^{-4}$ genau ein MeV ist; eine geometrische Deutung dafür steht aus. Die frühere Lesart
$E_0=1/\xi=7500$\,GeV verwechselt $E_0$ mit $E_0^2$ und setzt eine Einheit.
\end{abstract}

% ============================================================
\section*{Statusmarkierungen}
% ============================================================
\begin{center}
\begin{tabular}{@{}lp{11cm}@{}}
\textbf{[B]} & Algebraisch bewiesen. \\[2pt]
\textbf{[K]} & Numerisch bestätigt (Prüfskript, Vergleich mit Messwerten). \\[2pt]
\textbf{[S]} & Setzung, offen oder spekulativ. \\[2pt]
\textbf{[X]} & Widerlegt oder nicht tragfähig. \\[2pt]
\textbf{[Q]} & Konvention. \\
\end{tabular}
\end{center}
Konstanten: CODATA 2022. Prüfskript: \texttt{pruef\_386\_alpha\_geometrie.py}.

% ============================================================
\section{Die Grundform der FFGFT}
% ============================================================

Die FFGFT rechnet mit Verhältnissen. Masse, Energie, Temperatur und inverse Zeit sind in
natürlichen Einheiten eine Größe in vier Kleidungen (A085), und die Kopplung wird als Recheneinheit gesetzt:
$\alpha=1$ \QM\ (A267). Der einzige Parameter ist $\xi=4/30000$. In dieser Grundform gibt es
keine fraktale Korrektur; $K_{\text{frak}}$, die SI-Werte von $\alpha$ und $G$ und die
ausgerollten Massentabellen gehören zur Übersetzung in SI, nicht zum Kern (A130). Die Frage
dieses Dokuments ist deshalb nicht, wie genau $1/137$ getroffen wird, sondern wo $\alpha$ bleibt,
wenn man es auf~1 setzt.

% ============================================================
\section{Wohin \texorpdfstring{$\alpha$}{alpha} bei der Umdefinition wandert}
% ============================================================

Die Feinstrukturkonstante ist
\begin{equation}
  \alpha=\frac{e^2}{4\pi\varepsilon_0\hbar c}.
\end{equation}
In natürlichen Einheiten mit $\hbar=c=\varepsilon_0=1$ (Heaviside-Lorentz) bleibt
\begin{equation}
  e^2=4\pi\alpha \quad\BM .
\end{equation}
Das $4\pi$ ist die volle Kugeloberfläche. Nach dem Gaußschen Satz ist die Ladung der Fluss durch
eine geschlossene Kugel; der Fluss verteilt sich auf den Raumwinkel $4\pi$. Setzt man $\alpha=1$,
wird
\begin{equation}
  e_{\text{geo}}=\sqrt{4\pi}=3{,}545 ,
\end{equation}
die Ladung ist dann geometrisch gemessen: ihr Quadrat ist der volle Raumwinkel. Mit dem
gemessenen $\alpha$ ist die historische Ladung $e_{\text{hist}}=\sqrt{4\pi\alpha}=0{,}3028$. Damit
\begin{equation}
  \alpha=\left(\frac{e_{\text{hist}}}{e_{\text{geo}}}\right)^2,\qquad
  \frac{e_{\text{geo}}}{e_{\text{hist}}}=\frac1{\sqrt\alpha}=11{,}706\quad\BM .
\end{equation}
Bei der Umdefinition verschwindet $\alpha$ also nicht. Es wandert in die Ladungseinheit: $\alpha$
ist das Quadrat des Verhältnisses von historischer zu geometrischer Ladungseinheit. In jeder
Rechnung bleibt es über $e^2=4\pi\alpha$ als Spur erhalten, pro Vertex ein Faktor $\sqrt\alpha$
(A267).

Dasselbe $4\pi$ steht im Higgs-Vakuum-Weg: Der Faktor $64\pi^4=16\pi^3\cdot4\pi$ der Vakuumformel
von 2025 ist das $4\pi$ aus $4\pi\varepsilon_0$ (Dok.~385).

% ============================================================
\section{Das geometrische Bild: zwei Kugeln um das Elektron}
% ============================================================

Um das Elektron lassen sich zwei Kugeln angeben. Auf der ersten ist die Coulomb-Energie gleich
der Ruheenergie, das ist der klassische Elektronenradius
\begin{equation}
  r_e=\frac{e^2}{4\pi\varepsilon_0\,m_ec^2}=2{,}818\times10^{-15}\ \text{m}.
\end{equation}
Auf der zweiten ist die Quantenenergie $\hbar c/r$ gleich der Ruheenergie, das ist die reduzierte
Compton-Wellenlänge
\begin{equation}
  \lambda_C=\frac{\hbar}{m_ec}=3{,}862\times10^{-13}\ \text{m}.
\end{equation}
Ihr Verhältnis ist genau die Feinstrukturkonstante:
\begin{equation}
  \frac{r_e}{\lambda_C}=\alpha ,\qquad \frac{a_0}{\lambda_C}=\frac1\alpha \quad\BM .
\end{equation}
Gleichwertig ist $\alpha$ das Verhältnis von Coulomb-Energie zu Quantenenergie auf derselben
Kugel, $\bigl(e^2/(4\pi\varepsilon_0r)\bigr)/(\hbar c/r)$, und zwar für jeden Radius~$r$ \KM.

$\alpha=1$ heißt geometrisch: Die Coulomb-Kugel und die Compton-Kugel fallen zusammen, und mit
ihnen der Bohr-Radius; $r_e$, $\lambda_C$ und $a_0$ sind dann eine einzige Skala $1/m_e$
(A267). Der SI-Wert $1/137$ sagt nur, dass in der historischen Ladungseinheit die Coulomb-Kugel
137-mal kleiner ist als die Compton-Kugel. Das ist die transparente geometrische Bedeutung von
$\alpha$.

% ============================================================
\section{Die FFGFT-Brücke in Verhältnisform}
% ============================================================

Die FFGFT verbindet $\alpha$ mit $\xi$ über die Energieskala
$E_0=\sqrt{m_em_\mu}$, das geometrische Mittel der beiden leichtesten geladenen Leptonen. In der
Grundform, also in T0-Einheiten mit $\alpha=1$ und ohne fraktale Korrektur, lautet die Brücke
\begin{equation}
  \xi\,E_0^2=1\quad\Longrightarrow\quad E_0^2=\frac1\xi=7500,\qquad E_0=86{,}60
  \quad\BM .
\end{equation}
Die Zahl $7500=1/\xi$ gehört also zu $E_0^2$, nicht zu $E_0$. Für sich genommen legt diese
Gleichung die T0-Energieeinheit fest: Eine Einheit ist
$E_0/86{,}60=\sqrt{\xi\,m_em_\mu}=84{,}85$\,keV \KM. Inhalt bekommt sie erst im Vergleich mit
einer zweiten Skala, und das ist der Übergang nach SI.

% ============================================================
\section{Die SI-Brücke und die Bezugsenergie}
% ============================================================

In SI wird aus der Verhältnisform (Dok.~011, A130)
\begin{equation}
  \alpha_{\text{SI}}=\xi\left(\frac{E_0}{E_{\text{ref}}}\right)^2,\qquad E_{\text{ref}}=1\ \text{MeV}.
\end{equation}
Hier wird die Brücke undurchsichtig. $\xi E_0^2$ hat die Dimension Energie$^2$, also braucht man
eine Bezugsenergie, und der Korpus nimmt dafür das MeV. Rechnet man rückwärts, welche
Bezugsenergie die Gleichung mit dem gemessenen $\alpha$ verlangt, findet man
\begin{equation}
  E_{\text{ref}}=E_0\sqrt{\xi/\alpha}=
  \begin{cases}
  0{,}99323\ \text{MeV} & \text{nackt, } E_0=7{,}348\ \text{MeV},\\
  0{,}99992\ \text{MeV} & \text{korrigiert, } E_0=7{,}397\ \text{MeV}
  \end{cases}\quad\KM .
\end{equation}
Der nackte Wert ist $\sqrt{K_{\text{frak}}}$\,MeV bis auf $8\times10^{-5}$. Die fraktale
Korrektur kommt also erst hier als Feinabgleich zwischen T0-Einheit und MeV hinzu; ohne sie gilt
$\alpha^{-1}=7500/E_0^2=138{,}9$, mit ihr $137{,}06$, $1{,}7\times10^{-4}$ neben dem Messwert.
Dasselbe in den Einheiten der Grundform: Die T0-Energieeinheit ist $\sqrt\alpha\cdot E_{\text{ref}}$.

Die Übereinstimmung von $\alpha$ sagt damit eines aus: Die Bezugsenergie ist auf $10^{-4}$ genau
ein MeV. Das ist eine Eigenschaft des Ankers $E_0$ (R135), keine unabhängige Vorhersage \KM. Eine
geometrische Bedeutung von $E_{\text{ref}}=1$\,MeV ist offen \SM. Naheliegende Kandidaten tragen
nicht: Die Paarschwelle $2m_e=1{,}022$\,MeV liegt $2{,}2\,\%$ daneben \XM. Die Aussage in
Dok.~013 und 014, $S_{T0}=1\,\text{MeV}/c^2$ sei vorhergesagt, ist zirkulär, weil
$m_e^{\text{T0}}=0{,}511$ schon in MeV angegeben ist; Dok.~014 trägt dazu einen Vermerk, Dok.~013
jetzt ebenfalls.

% ============================================================
\section{Folge für die Lesart \texorpdfstring{$E_0=1/\xi$}{E0 = 1/xi}}
% ============================================================

In Dok.~005 und 032 steht $E_0=1/\xi=7500$\,GeV; Dok.~133 und 165 haben die Gleichung
übernommen. Aus der Verhältnisform ergibt sich, was daran richtig ist und was nicht. Richtig ist
die Zahl: In T0-Einheiten ist $E_0^2=1/\xi=7500$ \BM. Falsch sind zwei Schritte \XM. Erstens
fehlt das Quadrat: Mit $E_0=1/\xi$ wäre $\xi E_0^2=7500$, weder~1 noch $\alpha$. Zweitens ist die
Einheit GeV an eine reine Zahl gehängt; $7500$\,GeV liegt um den Faktor $1{,}01\times10^6$ über
$E_0=7{,}397$\,MeV.

% ============================================================
\section{Was trägt und was offen ist}
% ============================================================

\begin{center}
\small
\begin{tabular}{@{}p{7.4cm}p{1.2cm}p{5.2cm}@{}}
\toprule
\textbf{Aussage} & \textbf{Status} & \textbf{Grund} \\
\midrule
$e^2=4\pi\alpha$; mit $\alpha=1$ ist $e=\sqrt{4\pi}$ & \BM & $\hbar=c=\varepsilon_0=1$ \\
$\alpha=(e_{\text{hist}}/e_{\text{geo}})^2$ & \BM & Umdefinition der Ladung \\
$\alpha=r_e/\lambda_C$, $a_0/\lambda_C=1/\alpha$ & \BM & zwei Kugeln um das Elektron \\
$\alpha=1$ als Recheneinheit & \QM & A267 \\
$\xi E_0^2=1$, $E_0^2=1/\xi=7500$ in T0-Einheiten & \BM & Grundform \\
T0-Energieeinheit $\sqrt{\xi m_em_\mu}=84{,}85$\,keV & \KM & Prüfskript \\
$E_{\text{ref}}=0{,}99992$\,MeV (korrigiert), $0{,}99323$\,MeV (nackt) & \KM & Eigenschaft des Ankers (R135) \\
geometrische Bedeutung von $E_{\text{ref}}=1$\,MeV & \SM & offen \\
$E_{\text{ref}}=2m_e$ & \XM & $2{,}2\,\%$ daneben \\
$S_{T0}=1\,\text{MeV}/c^2$ vorhergesagt & \XM & zirkulär \\
$E_0=1/\xi=7500$\,GeV & \XM & Quadrat fehlt, Einheit gesetzt \\
\bottomrule
\end{tabular}
\end{center}

% ============================================================
\section{Angeglichene Stellen im Korpus}
% ============================================================

Mit diesem Dokument tragen die Stellen zu $E_0=1/\xi$ einen datierten Hinweis auf die
Verhältnisform $E_0^2=1/\xi$: Dok.~005 (bisher ohne Vermerk), 032, 133 und 165, jeweils De/En.
Dok.~384 stellt die Grundform voran. Dok.~013 vermerkt, dass $S_{T0}=1\,\text{MeV}/c^2$ per
Konstruktion gilt.

% ============================================================
\section{Fazit}
% ============================================================

$\alpha$ verschwindet bei der Umdefinition nicht, es wird zur Ladungseinheit. Mit $\alpha=1$ ist
die Ladung über die volle Kugeloberfläche gemessen, $e=\sqrt{4\pi}$, und die Coulomb-Kugel des
Elektrons fällt mit seiner Compton-Kugel zusammen. Der Messwert $1/137$ ist das
Radienverhältnis $r_e/\lambda_C$ in der historischen Ladungseinheit. Die FFGFT-Brücke lautet in
ihrer Grundform $\xi E_0^2=1$, also $E_0^2=1/\xi$; das ist transparent. Undurchsichtig ist nur der
Schritt nach SI, weil er eine Bezugsenergie braucht, die auf $10^{-4}$ genau ein MeV ist, ohne
dass dafür eine geometrische Begründung vorliegt.

\vspace{1em}
\noindent\textit{Prüfskript:} \texttt{2/python/Dok386\_Skripte/pruef\_386\_alpha\_geometrie.py}
--- 19/19.
